\section*{Projeler}

\datedsubsection{Otonom Sualtı Aracı (AUV)}{2025}
\textit{Teknolojiler: Python}

TEKNOFEST İnsansız Su Altı Sistemleri Yarışması için planlanan projede sistem tasarımı, kontrol mimarisi ve Python tabanlı yazılım yaklaşımı üzerine çalıştım. Çalışma araştırma ve planlama aşamasında kaldı; fiziksel araç geliştirme veya saha testi gerçekleştirilmedi.

\vspace{0.2cm}
\noindent \textbf{Odak Noktaları:}
\begin{itemize}
    \item Kontrol mimarisi ve yazılım yaklaşımının planlanması
    \item Sensör ve donanım seçeneklerinin araştırılması
    \item Yarışma süreci için sistem tasarımı ve teknik dokümantasyon
\end{itemize}

\vspace{0.5cm}

\datedsubsection{Dijital Diş Dünyası}{2024}
\textit{Teknolojiler: Vue.js, Laravel, MySQL}

Dijital Diş Dünyası, diş hekimliği sektörüne yönelik aktif bir web platformudur. Vue.js ön yüz, Laravel arka uç ve MySQL veritabanı bileşenlerini full-stack geliştirme kapsamında geliştirdim ve entegre ettim. Aktif site: \href{https://dijitaldisdunyasi.com/}{dijitaldisdunyasi.com}.

\vspace{0.2cm}
\noindent \textbf{Odak Noktaları:}
\begin{itemize}
    \item Vue.js ve Laravel ile full stack geliştirme
    \item Ön yüz, arka uç ve veritabanı entegrasyonu
    \item Aktif yayında olan web platformu deneyimi
\end{itemize}
